\documentclass[11pt,a4paper]{article}

%  XeLaTeX
\usepackage{fontspec}

\usepackage{algorithm}
\usepackage{algorithmic}
\usepackage{amsmath,amssymb,amsthm}
\usepackage{bm}
\usepackage[useregional]{datetime2}
\usepackage{enumitem}
\usepackage{booktabs}
\usepackage{float}
\usepackage{graphicx}
\usepackage{tabularray}
\usepackage{framed}
\usepackage[a4paper, margin=2.5cm]{geometry}
\usepackage{eso-pic}
\usepackage{mathtools}
\usepackage{tikz}
\usetikzlibrary{calc}
\usepackage{xcolor}

\usepackage[backend=biber, style=apa, sorting=nyt]{biblatex}
\addbibresource{D:/github/ENEL445/report/references.bib}
\graphicspath{{D:/github/ENEL445/figs/}}

\usepackage[colorlinks=true, linkcolor=blue, citecolor=blue, urlcolor=blue]{hyperref}
\hypersetup{
  pdfauthor   = {David Ewing},
  pdfsubject  = {\jobname.tex},
  pdfkeywords = {source: \jobname.tex}
}

% Running margin label
\newcommand{\mymarginlabel}{David Ewing dew59@uclive.ac.nz | \DTMnow\ | \jobname.tex}
\AddToShipoutPictureBG{%
  \begin{tikzpicture}[remember picture,overlay]
    \node[rotate=90, anchor=north]
      at ($(current page.north west)+(1.4cm,-13cm)$)
      {\scriptsize\ttfamily \mymarginlabel};
  \end{tikzpicture}%
}

% Custom commands
\newcommand{\E}{\mathbb{E}}
\newcommand{\R}{\mathbb{R}}
\newcommand{\N}{\mathcal{N}}
\newcommand{\ELBO}{\text{ELBO}}
\newcommand{\tr}{\text{tr}}
\newcommand{\sigmoid}{\sigma}
\setlength{\parindent}{0pt}

\title{
Case Study 5 --- Hierarchical Bayesian Logistic Regression:\\[0.4em]
Gradient-Based Optimisation and Posterior Variance Collapse
}
\author{
  David Ewing (82171165)
}
\date{11 May 2026}

\begin{document}

\maketitle

\begin{abstract}
\setlength{\parindent}{0pt}
\setlength{\parskip}{0.7em}

This report applies the gradient-based optimisation methods from ENEL445 to variational
inference in hierarchical Bayesian logistic regression.
This model combines two sources of analytical difficulty from previous case studies:
the non-conjugate logistic likelihood (Case Study~3) and the hierarchical random-effects
structure (Case Study~4).
Group random effects $u_j \sim \mathcal{N}(0,\tau_u^{-1})$ augment the fixed effects
$\bm{\beta}$, and the Jaakkola--Jordan \parencite{Jaakkola2000} bound is required to
retain a tractable ELBO.
The unconstrained parameter dimension is $D=17$.

Four optimisation methods are applied: CAVI, gradient ascent with Armijo backtracking,
Newton's method, and BFGS.
The reference posterior is provided by a P\'{o}lya--Gamma blocked Gibbs sampler
\parencite{Polson2013} with $T=50$ truncation terms, 3 chains $\times$ 2000 iterations.

A key finding is pronounced posterior variance collapse in both $\beta_0$ and $\beta_1$
(SD ratios $\approx 0.46$), reflecting the mean-field VI family's failure to capture
the posterior correlation between the fixed effects and the group random effects.
CAVI and BFGS converge to the same ELBO optimum $1050\times$ and $81\times$ faster
than the PG Gibbs sampler respectively.
Newton's method again fails to find the correct optimum, converging to a suboptimal
point with markedly worse ELBO.

\end{abstract}

\newpage

%─────────────────────────────────────────────────────────────────────────────
\section*{Introduction}
%─────────────────────────────────────────────────────────────────────────────

Hierarchical logistic regression is a natural model for binary outcomes observed in
grouped or clustered data: students within schools, patients within hospitals, or
binary responses from repeated measurements on the same subject.
Each group $j$ has a random intercept $u_j$ that shifts the group-specific log-odds
relative to the population, while $\bm{\beta}$ captures the population-level
covariate effect.

Bayesian inference in this model requires integrating over both $\bm{\beta}$ and
the random effects $\mathbf{u} = (u_1,\ldots,u_J)^T$, and over the random-effect
precision $\tau_u$.
The logistic likelihood is not conjugate to any natural prior, so both MCMC and
variational approximations are computationally more demanding than the Gaussian case.

This report is Case Study~5 in this project, combining the Jaakkola--Jordan bound
from Case Study~3 with the hierarchical structure from Case Study~4.
The mean-field ELBO is maximised using four optimisation strategies, and the results
are compared against the P\'{o}lya--Gamma Gibbs reference.

%─────────────────────────────────────────────────────────────────────────────
\section{Model Formulation}
%─────────────────────────────────────────────────────────────────────────────

\subsection{Hierarchical Bayesian Logistic Regression}

The model is
\begin{align}
  y_{ij} \mid \bm{\beta}, u_j
    &\sim \mathrm{Bernoulli}\!\bigl(\sigmoid(\mathbf{x}_{ij}^T\bm{\beta} + u_j)\bigr),
    \quad i=1,\ldots,n_j,\; j=1,\ldots,J,\\
  u_j \mid \tau_u &\sim \mathcal{N}(0,\; \tau_u^{-1}),\\
  \bm{\beta} &\sim \N(\bm{0},\; \lambda^{-1}\bm{I}_p), \quad \lambda=0.1,\\
  \tau_u &\sim \mathrm{Gamma}(\alpha_u,\; \gamma_u), \quad \alpha_u=1, \;\gamma_u=0.1,
\end{align}
where $\sigmoid(\psi)=(1+e^{-\psi})^{-1}$ is the logistic sigmoid and
$\mathbf{x}_{ij} = (1, x_{ij})^T \in \R^2$.

The test case uses $N=150$ binary observations in $J=5$ groups (30 per group),
$x_{ij} \sim \mathrm{Uniform}(-2,2)$, and true parameters
$\bm{\beta}_{\mathrm{true}} = (-1.0,\; 2.0)$,
$\tau_{u,\mathrm{true}}=2.0$ ($\sigma_u \approx 0.707$).

\subsection{Mean-Field Variational Approximation}

The mean-field variational family is
\begin{align}
  q(\bm{\beta},\mathbf{u},\tau_u)
  &= q(\bm{\beta})\;\prod_{j=1}^{J} q(u_j)\;\cdot\; q(\tau_u),
\end{align}
with
\begin{align}
  q(\bm{\beta}) &= \N(\bm{\mu}_\beta,\,\bm{\Sigma}_\beta), \quad
    \bm{\Sigma}_\beta = \bm{L}\bm{L}^T,\\
  q(u_j) &= \mathcal{N}(m_j,\, v_j),\\
  q(\tau_u) &= \mathrm{Gamma}(a_u,\, b_u).
\end{align}

\subsection{Jaakkola--Jordan ELBO}

The log-likelihood $\E_q[\log\sigmoid(\psi_{ij})]$ where
$\psi_{ij} = \mathbf{x}_{ij}^T\bm{\beta} + u_j$ is intractable under a Gaussian $q$.
The Jaakkola--Jordan bound introduces local variational parameters $\xi_{ij}>0$:
\begin{equation}
  \log\sigmoid(\psi_{ij}) \geq \tfrac{\psi_{ij}}{2}
    - \log 2\cosh\!\tfrac{\xi_{ij}}{2}
    - \tfrac{\xi_{ij}}{2}\bigl(\E_q[\psi_{ij}^2]/\xi_{ij}^2 - 1\bigr)\tfrac{\tanh(\xi_{ij}/2)}{4}.
\end{equation}
At the optimal $\xi_{ij}^* = \sqrt{\E_q[\psi_{ij}^2]}$, where
\begin{equation}
  \E_q[\psi_{ij}^2]
  = (\mathbf{x}_{ij}^T\bm{\mu}_\beta + m_j)^2
    + \mathbf{x}_{ij}^T\bm{\Sigma}_\beta\mathbf{x}_{ij} + v_j,
\end{equation}
the ELBO is
\begin{align}
  \mathcal{L}(\bm{\theta})
  &= \sum_{ij}\!\bigl[(y_{ij}-\tfrac{1}{2})\bar{\psi}_{ij}
    - \log 2 - \log\cosh\tfrac{\xi_{ij}^*}{2}\bigr]
  - \mathrm{KL}(q(\bm{\beta})\|p(\bm{\beta}))\nonumber\\
  &\quad + \E_q[\log p(\mathbf{u}\mid\tau_u)]
    + \E_q[\log p(\tau_u)]
    + \sum_j H[q(u_j)] + H[q(\tau_u)],
\end{align}
where $\bar{\psi}_{ij} = \mathbf{x}_{ij}^T\bm{\mu}_\beta + m_j$ and
all terms are available in closed form.

\subsection{CAVI Update Equations}

Let $\lambda_{ij} = \tanh(\xi_{ij}^*/2)/(4\xi_{ij}^*)$.
The coordinate ascent updates are:
\begin{align}
  \bm{\Sigma}_\beta^{-1}
    &\leftarrow \lambda\bm{I} + 2\sum_{ij}\lambda_{ij}\,\mathbf{x}_{ij}\mathbf{x}_{ij}^T,
    \label{eq:cavi_Sigma}\\
  \bm{\mu}_\beta
    &\leftarrow \bm{\Sigma}_\beta\!\sum_{ij}
      (y_{ij}-\tfrac{1}{2} - 2\lambda_{ij}m_j)\,\mathbf{x}_{ij},\\
  v_j^{-1}
    &\leftarrow 2\sum_{i\in j}\lambda_{ij} + \E[\tau_u], \quad
  m_j \leftarrow v_j\!\sum_{i\in j}
    (y_{ij}-\tfrac{1}{2} - 2\lambda_{ij}\mathbf{x}_{ij}^T\bm{\mu}_\beta),\\
  a_u &\leftarrow \alpha_u + \tfrac{J}{2}, \quad
  b_u \leftarrow \gamma_u + \tfrac{1}{2}\sum_j (m_j^2 + v_j).
\end{align}

\subsection{Unconstrained Parameterisation}

The unconstrained vector $\bm{\theta} \in \R^{17}$ is
\begin{equation}
  \bm{\theta} = \bigl(\bm{\mu}_\beta,\;
    \tilde{\ell}_{11}, \ell_{21}, \tilde{\ell}_{22},\;
    m_1,\ldots,m_5,\;
    \log v_1,\ldots,\log v_5,\;
    \log a_u, \log b_u
  \bigr),
\end{equation}
giving $D = 2 + 3 + 5 + 5 + 2 = 17$.

\subsection{Key Notation}

\begin{center}
\begin{tabular}{@{}ll@{}}
\toprule
Symbol & Meaning \\
\midrule
\multicolumn{2}{@{}l}{\textbf{Bayesian model}} \\
$\bm{X}$ & Design matrix ($N \times p$), rows $\mathbf{x}_{ij}^T$ \\
$\mathbf{y}$ & Binary response vector ($y_{ij} \in \{0,1\}$) \\
$\bm{\beta}$ & Fixed-effect coefficient vector \\
$u_j$ & Group random effect for group $j$ \\
$\tau_u$ & Random-effect precision \\
$\sigmoid(\cdot)$ & Logistic sigmoid $\sigmoid(\psi)=(1+e^{-\psi})^{-1}$ \\
$N, J, p$ & Total observations, groups, predictors \\
\midrule
\multicolumn{2}{@{}l}{\textbf{Variational approximation}} \\
$q(\bm{\beta})$ & Variational Gaussian for fixed effects \\
$\bm{\mu}_\beta, \bm{\Sigma}_\beta$ & Mean and covariance of $q(\bm{\beta})$ \\
$\bm{L}$ & Cholesky factor, $\bm{\Sigma}_\beta=\bm{L}\bm{L}^T$ \\
$m_j, v_j$ & Mean and variance of $q(u_j)$ \\
$a_u, b_u$ & Gamma parameters of $q(\tau_u)$ \\
$\xi_{ij}$ & Local JJ variational parameter; optimal at $\sqrt{\E_q[\psi_{ij}^2]}$ \\
$\lambda_{ij}$ & JJ weight $\tanh(\xi_{ij}/2)/(4\xi_{ij})$ \\
$\mathcal{L}(\bm{\theta})$ & Jaakkola--Jordan ELBO (maximised) \\
\midrule
\multicolumn{2}{@{}l}{\textbf{Optimisation}} \\
$\bm{\theta}$ & Unconstrained vector ($D=17$) \\
$\alpha_t$ & Step size at iteration $t$ \\
\midrule
\multicolumn{2}{@{}l}{\textbf{P\'{o}lya--Gamma sampler}} \\
$\omega_{ij}$ & PG latent variable: $\omega_{ij}\mid\bm{\beta},u_j \sim \mathrm{PG}(1,\psi_{ij})$ \\
$\kappa_{ij}$ & Pseudo-response $\kappa_{ij} = y_{ij} - 1/2$ \\
$T$ & Truncation order for PG series ($T=50$) \\
\midrule
\multicolumn{2}{@{}l}{\textbf{Lecture correspondence} (Prof.\ R.\,S.\ Smith, ENEL445 2026)} \\
$J(\theta)$ & Cost function $=\mathcal{L}(\bm{\theta})$ in this report \\
$\Phi$ & Regressor matrix $=\bm{X}$ in this report \\
\bottomrule
\end{tabular}
\end{center}

%─────────────────────────────────────────────────────────────────────────────
\section{P\'{o}lya--Gamma Blocked Gibbs Sampler}
%─────────────────────────────────────────────────────────────────────────────

The P\'{o}lya--Gamma (PG) identity of \textcite{Polson2013},
\begin{equation}
  \sigmoid(\psi)^{y}(1-\sigmoid(\psi))^{1-y}
  \propto e^{\kappa\psi}
    \int_0^\infty e^{-\omega\psi^2/2}\,p(\omega)\,\mathrm{d}\omega,
  \quad \omega \sim \mathrm{PG}(1,0),
\end{equation}
introduces a latent variable $\omega_{ij}$ that renders the logistic likelihood
conditionally Gaussian.
Each Gibbs sweep samples:
\begin{align}
  \omega_{ij} \mid \bm{\beta}, u_j
    &\sim \mathrm{PG}(1,\; \mathbf{x}_{ij}^T\bm{\beta} + u_j), \label{eq:pg_omega}\\
  \bm{\beta} \mid \bm{\omega}, \mathbf{u}, \mathbf{y}
    &\sim \mathcal{N}(V_\beta(\bm{X}^T\bm{\kappa} - \bm{X}^T\mathrm{diag}(\bm{\omega})\mathbf{u}[\mathrm{g}]),\;V_\beta),
    \quad V_\beta = (\lambda\bm{I} + \bm{X}^T\mathrm{diag}(\bm{\omega})\bm{X})^{-1},\\
  u_j \mid \bm{\beta}, \bm{\omega}_j, \tau_u, \mathbf{y}
    &\sim \mathcal{N}(V_j^*m_j^*,\; V_j^*),
    \quad V_j^* = \bigl(\textstyle\sum_{i\in j}\omega_{ij} + \tau_u\bigr)^{-1},
    \quad m_j^* = \textstyle\sum_{i\in j}(\kappa_{ij} - \omega_{ij}\mathbf{x}_{ij}^T\bm{\beta}),\\
  \tau_u \mid \mathbf{u}
    &\sim \mathrm{Gamma}\!\left(\alpha_u + \tfrac{J}{2},\;
      \gamma_u + \tfrac{1}{2}\textstyle\sum_j u_j^2\right).
\end{align}

PG draws use the truncated series with $T=50$ terms.
Three chains run for 2000 iterations each with 500 burn-in, yielding
$3 \times 1500 = 4500$ post-burnin samples.

%─────────────────────────────────────────────────────────────────────────────
\section{Optimisation Problem Formulation}
%─────────────────────────────────────────────────────────────────────────────

\subsection*{Decision Variables and Constraints}

Optimisation is over $\bm{\theta} \in \R^{17}$.
All constraints ($\bm{\Sigma}_\beta \succ 0$, $v_j > 0$, $a_u, b_u > 0$) are satisfied
automatically by the parameterisation.

\subsection*{Objective Function}

Maximise the Jaakkola--Jordan ELBO $\mathcal{L}(\bm{\theta}):\R^{17}\to\R$ at optimal
$\xi_{ij}^* = \sqrt{\bar{\psi}_{ij}^2 + \mathbf{x}_{ij}^T\bm{\Sigma}_\beta\mathbf{x}_{ij} + v_j}$.

\subsection*{Optimisation Components}

\begin{description}
\item[CAVI.]
  Closed-form coordinate updates (Equations~\ref{eq:cavi_Sigma}--above).
  Unit step; exit when $|\Delta\mathcal{L}| < 10^{-8}$.

\item[Gradient ascent.]
  Direction: central FD gradient ($h=10^{-5}$).
  Armijo backtracking, $\alpha_0=0.5$, $c_1=0.1$, $\rho=0.5$.
  Exit: $|\Delta\mathcal{L}| < 10^{-7}$ or 2000 iterations.

\item[Newton's method.]
  FD Hessian ($17\times17$); regularised to ensure ascent direction.
  Armijo backtracking; exit at 50 iterations.

\item[BFGS.]
  Inverse-Hessian approximation via rank-2 BFGS update.
  Wolfe conditions (scipy); exit at $\|\bm{g}\|<10^{-7}$ or 2000 iterations.
\end{description}

%─────────────────────────────────────────────────────────────────────────────
\section{Numerical Study}
%─────────────────────────────────────────────────────────────────────────────

\subsection{Test Case Design}

\begin{itemize}
  \item Model: $y_{ij} \sim \mathrm{Bernoulli}(\sigmoid(\beta_0 + \beta_1 x_{ij} + u_j))$
  \item True parameters: $\bm{\beta}_{\mathrm{true}} = (-1.0,\; 2.0)$,
    $\tau_{u,\mathrm{true}}=2.0$
  \item Groups: $J=5$, $n_j=30$, $N=150$
  \item Covariates: $x_{ij} \sim \mathrm{Uniform}(-2, 2)$
  \item Prior: $\bm{\beta} \sim \N(\bm{0}, 10\bm{I})$;
    $\tau_u \sim \mathrm{Gamma}(1, 0.1)$
  \item Seed: \texttt{82171165} (student ID)
\end{itemize}

\subsection{Gradient Verification}

The gradient of the ELBO with respect to the 17-dimensional $\bm{\theta}$ was verified
against central finite differences at a CAVI warm-start point.
All 17 components agreed to relative error below $10^{-3}$.

\begin{figure}[H]
  \centering
  \includegraphics[width=0.95\linewidth]{hierarchical_logistic_elbo_gradient_check.png}
  \caption{Finite-difference gradient verification at the CAVI warm-start:
    $h=10^{-5}$ versus $h=10^{-7}$.
    All 17 components agree, confirming correct gradient implementation.}
  \label{fig:gradcheck}
\end{figure}

\subsection{ELBO Landscape}

Figure~\ref{fig:landscape} shows the 2-D ELBO slice in the
$(\mu_{\beta_0}, \mu_{\beta_1})$ plane with all other parameters at the warm start.

\begin{figure}[H]
  \centering
  \includegraphics[width=0.62\linewidth]{hierarchical_logistic_elbo_landscape.png}
  \caption{ELBO landscape in the $(\mu_{\beta_0},\mu_{\beta_1})$ plane.
    Red star: CAVI warm start; white cross: true $\bm{\beta}$.}
  \label{fig:landscape}
\end{figure}

%─────────────────────────────────────────────────────────────────────────────
\section{Results}
%─────────────────────────────────────────────────────────────────────────────

\subsection{Generated Data}

Figure~\ref{fig:data} shows the 150 binary observations grouped by $j$.
Group random effects $u_j$ shift the group-specific logistic curve relative to
the population curve $\sigmoid(\beta_0 + \beta_1 x)$.

\begin{figure}[H]
  \centering
  \includegraphics[width=0.80\linewidth]{hierarchical_logistic_data_scatter.png}
  \caption{Generated hierarchical logistic data ($N=150$, $J=5$).
    Dotted lines show group-specific sigmoid curves; solid: population mean.}
  \label{fig:data}
\end{figure}

\subsection{ELBO Convergence}

Figure~\ref{fig:elbocv} shows the ELBO history for all methods.
CAVI converges in 37 iterations; BFGS in 31.
Both reach ELBO $= -54.91$.
Newton's method converges in 9 iterations but to a suboptimal point
(ELBO $= -77.36$) --- the same failure mode seen in Case Study~4.
Gradient ascent requires 416 iterations to reach the same optimum as CAVI.

\begin{figure}[H]
  \centering
  \includegraphics[width=\linewidth]{hierarchical_logistic_elbo_convergence.png}
  \caption{ELBO convergence: CAVI (left) and gradient-based methods (right).
    Newton converges to a suboptimal point.}
  \label{fig:elbocv}
\end{figure}

\subsection{Step-Size and Condition-Number Behaviour}

\begin{figure}[H]
  \centering
  \includegraphics[width=0.75\linewidth]{hierarchical_logistic_gradient_stepsize.png}
  \caption{Armijo step sizes for gradient ascent (log scale).}
  \label{fig:stepsize}
\end{figure}

\begin{figure}[H]
  \centering
  \includegraphics[width=0.65\linewidth]{hierarchical_logistic_newton_condition.png}
  \caption{Hessian condition number at each Newton iteration.
    Poor conditioning explains Newton's convergence to a suboptimal point.}
  \label{fig:cond}
\end{figure}

\subsection{P\'{o}lya--Gamma Gibbs Sampler Diagnostics}

Figures~\ref{fig:trace_b0}--\ref{fig:trace_tauu} show trace plots for three
monitored parameters.
ACF plots (Figures~\ref{fig:acf_b0}--\ref{fig:acf_tauu}) show rapid decorrelation.
Figure~\ref{fig:rhat} confirms $\hat{R} < 1.01$ for all parameters.

\begin{figure}[H]
  \centering
  \includegraphics[width=\linewidth]{hierarchical_logistic_gibbs_trace_beta0.png}
  \caption{PG Gibbs trace: $\beta_0$. Red dashed: true value.}
  \label{fig:trace_b0}
\end{figure}

\begin{figure}[H]
  \centering
  \includegraphics[width=\linewidth]{hierarchical_logistic_gibbs_trace_beta1.png}
  \caption{PG Gibbs trace: $\beta_1$.}
  \label{fig:trace_b1}
\end{figure}

\begin{figure}[H]
  \centering
  \includegraphics[width=\linewidth]{hierarchical_logistic_gibbs_trace_tau_u.png}
  \caption{PG Gibbs trace: $\tau_u$.}
  \label{fig:trace_tauu}
\end{figure}

\begin{figure}[H]
  \centering
  \includegraphics[width=0.95\linewidth]{hierarchical_logistic_gibbs_acf_beta0.png}
  \caption{ACF for $\beta_0$ (chain 1). Red dashed: 95\% band.}
  \label{fig:acf_b0}
\end{figure}

\begin{figure}[H]
  \centering
  \includegraphics[width=0.95\linewidth]{hierarchical_logistic_gibbs_acf_beta1.png}
  \caption{ACF for $\beta_1$.}
  \label{fig:acf_b1}
\end{figure}

\begin{figure}[H]
  \centering
  \includegraphics[width=0.95\linewidth]{hierarchical_logistic_gibbs_acf_tau_u.png}
  \caption{ACF for $\tau_u$.}
  \label{fig:acf_tauu}
\end{figure}

\begin{figure}[H]
  \centering
  \includegraphics[width=0.55\linewidth]{hierarchical_logistic_gibbs_rhat.png}
  \caption{Gelman--Rubin $\hat{R}$ diagnostic.
    All values below 1.01 confirm chain convergence.}
  \label{fig:rhat}
\end{figure}

\subsection{Random Effects Recovery}

Figure~\ref{fig:re} compares CAVI and Gibbs estimates of the five group random effects.
Both methods capture the direction and magnitude of each $u_j$; CAVI shows tighter
credible intervals consistent with the general variance collapse finding.

\begin{figure}[H]
  \centering
  \includegraphics[width=\linewidth]{hierarchical_logistic_random_effects.png}
  \caption{Random effect estimates $u_j$ for CAVI (left) and PG Gibbs (right).
    Stars: true values; bars: $\pm 2$ posterior SD.}
  \label{fig:re}
\end{figure}

\subsection{Posterior Comparison: VI versus Gibbs}

Figures~\ref{fig:vbgibbs_b0}--\ref{fig:vbgibbs_tauu} overlay the VI Gaussian
approximations against the PG Gibbs histograms.

\begin{figure}[H]
  \centering
  \includegraphics[width=0.80\linewidth]{hierarchical_logistic_vb_gibbs_beta0.png}
  \caption{Posterior comparison for $\beta_0$: Gibbs histogram versus VI.
    The VI curves are noticeably narrower than Gibbs (SD ratio $\approx 0.46$).}
  \label{fig:vbgibbs_b0}
\end{figure}

\begin{figure}[H]
  \centering
  \includegraphics[width=0.80\linewidth]{hierarchical_logistic_vb_gibbs_beta1.png}
  \caption{Posterior comparison for $\beta_1$. Similar collapse to $\beta_0$.}
  \label{fig:vbgibbs_b1}
\end{figure}

\begin{figure}[H]
  \centering
  \includegraphics[width=0.80\linewidth]{hierarchical_logistic_vb_gibbs_tau_u.png}
  \caption{Posterior comparison for $\tau_u$.
    CAVI and BFGS overestimate $\E[\tau_u]$ relative to Gibbs, possibly due to
    the variance collapse in $u_j$ causing an artificially small $\sum_j(m_j^2+v_j)$.}
  \label{fig:vbgibbs_tauu}
\end{figure}

\subsection{Posterior Standard-Deviation Ratios}

Figure~\ref{fig:sdratios} shows $\sigma_{\mathrm{VI}}/\sigma_{\mathrm{Gibbs}}$ for
each converged method.
CAVI, gradient ascent, and BFGS all give SD ratios $\approx 0.46$ for both $\beta_0$
and $\beta_1$, indicating that the VI posteriors are roughly half as wide as the
Gibbs reference.
This is a more symmetric collapse than in Case Study~4 (where only $\beta_0$ collapsed),
because both fixed effects are confounded with the group random effects in the
non-conjugate hierarchical logistic model.

\begin{figure}[H]
  \centering
  \includegraphics[width=0.85\linewidth]{hierarchical_logistic_sd_ratios.png}
  \caption{Posterior SD ratios (VI / Gibbs). Dashed: ideal 1.0.
    All converged VI methods show $\approx 0.46$ collapse for both fixed effects.}
  \label{fig:sdratios}
\end{figure}

\subsection{Posterior Mean Errors}

Figure~\ref{fig:meanerr} shows absolute mean errors.
CAVI, gradient ascent, and BFGS recover posterior means close to the Gibbs reference.
Newton's method has large mean errors due to its convergence to the wrong optimum.

\begin{figure}[H]
  \centering
  \includegraphics[width=0.85\linewidth]{hierarchical_logistic_mean_errors.png}
  \caption{Absolute posterior mean errors relative to PG Gibbs.}
  \label{fig:meanerr}
\end{figure}

\subsection{Performance Summary}

Table~\ref{tab:perf} summarises convergence and runtime.
Table~\ref{tab:post} summarises posterior means and SD ratios.

\begin{table}[H]
  \centering
  \caption{Optimisation performance: hierarchical logistic regression.}
  \label{tab:perf}
  \begin{tblr}{colspec={lrrrr}, hline{1,2,Z}={solid}}
  Method & Iterations & Runtime (ms) & Speedup vs Gibbs & Final ELBO \\
  CAVI & 37 & 13.4 & 1013x & -54.907 \\
  Gradient Ascent & 416 & 1576.6 & 9x & -54.907 \\
  Newton's Method & 9 & 535.3 & 25.3x & -77.360 \\
  BFGS & 31 & 163.3 & 83x & -54.907 \\
  PG Gibbs (3 chains) & 3x2000 & 13548.4 & 1x (reference) & N/A \\
\end{tblr}

\end{table}

\begin{table}[H]
  \centering
  \caption{Posterior summary: hierarchical logistic regression.}
  \label{tab:post}
  \begin{tblr}{colspec={lrrrrr}, hline{1,2,Z}={solid}}
  Method & $\mu_{\beta_0}$ & $\mu_{\beta_1}$ & $\mathrm{E}[\tau_u]$ & SD ratio $\beta_0$ & SD ratio $\beta_1$ \\
  CAVI & -1.4457 & 2.2019 & 10.0547 & 0.463 & 0.459 \\
  Gradient Ascent & -1.4457 & 2.2021 & 10.0298 & 0.463 & 0.459 \\
  Newton's Method & -0.7544 & 3.4213 & 1.7085 & 0.108 & 0.857 \\
  BFGS & -1.4458 & 2.2020 & 10.0545 & 0.463 & 0.459 \\
  PG Gibbs (ref) & -1.5484 & 2.3748 & 8.0822 & 1.000 & 1.000 \\
\end{tblr}

\end{table}

\subsection{Computational Cost}

Figure~\ref{fig:timing} shows the runtime comparison.
CAVI is the fastest at 15\,ms, $1050\times$ faster than the PG Gibbs sampler (15.8\,s).
BFGS is $81\times$ faster than Gibbs.
Gradient ascent is only $9\times$ faster despite finding the same optimum as CAVI.
Newton's method achieves $23\times$ speedup but converges to a suboptimal point.
The PG Gibbs is substantially more expensive than the conjugate Gibbs in Case Study~4
because each iteration requires $N \times T = 150 \times 50 = 7500$ Gamma variates
for the PG augmentation.

\begin{figure}[H]
  \centering
  \includegraphics[width=0.85\linewidth]{hierarchical_logistic_timing_comparison.png}
  \caption{Runtime comparison (log scale). Annotated speedup factors relative to PG Gibbs.}
  \label{fig:timing}
\end{figure}

%─────────────────────────────────────────────────────────────────────────────
\section{Discussion}
%─────────────────────────────────────────────────────────────────────────────

\paragraph{Posterior variance collapse.}
Both $\beta_0$ and $\beta_1$ exhibit SD ratios $\approx 0.46$ for all well-converged VI methods.
Unlike Case Study~4 (where collapse was confined to $\beta_0$), here \emph{both} fixed
effects are affected.
The mechanism is the same: mean-field VI discards the posterior correlation between
$\bm{\beta}$ and $\mathbf{u}$; in the logistic model, the JJ variational parameters
$\lambda_{ij}$ depend on $\bm{\Sigma}_\beta$ as well as $v_j$, so tightening the bound
in one factor indirectly narrows both.
In comparison, the simple logistic case (Case Study~3) showed near-zero collapse, because
without group structure there is no confounding between $\bm{\beta}$ and latent effects.

\paragraph{Overestimation of $\tau_u$.}
The converged VI methods give $\E[\tau_u] \approx 10$, roughly $25\%$ higher than the
Gibbs reference ($\approx 8$).
When the mean-field approximation collapses $q(u_j)$ to be too narrow
($v_j$ too small), the CAVI update for $b_u$ underestimates
$\sum_j(m_j^2 + v_j)$, driving $b_u$ down and $\E[\tau_u] = a_u/b_u$ up.
This is a secondary bias induced by the primary variance collapse.

\paragraph{Newton's method.}
Newton's method fails in the same manner as in Case Study~4: the $17\times17$
FD Hessian is poorly conditioned throughout, because the ELBO mixes parameters that
operate on very different scales.
The 9-iteration convergence to ELBO $= -77.36$ is fast but to the wrong point.

\paragraph{Comparison with Case Study~4.}
The hierarchical logistic model is computationally more demanding (PG augmentation
vs conjugate Gibbs), but the VI methods are similarly fast in absolute terms.
The qualitative finding --- CAVI/BFGS converge; Newton fails; gradient ascent is
much slower than BFGS --- is consistent across both hierarchical case studies.
The difference in variance collapse (symmetric in CS5, asymmetric in CS4) reflects the
differing roles of the intercept in linear versus logistic hierarchical models.

%─────────────────────────────────────────────────────────────────────────────
\section{Conclusion}
%─────────────────────────────────────────────────────────────────────────────

Hierarchical Bayesian logistic regression with $J=5$ groups and $D=17$ unconstrained
parameters confirms and extends the findings from Case Study~4.
CAVI and BFGS both find the mean-field ELBO optimum with speedups of
$1050\times$ and $81\times$ over the PG Gibbs sampler respectively.
Both converged VI methods exhibit posterior variance collapse for all fixed effects
(SD ratio $\approx 0.46$), more severe than the simple logistic case (Case Study~3)
and more symmetric than the hierarchical linear case (Case Study~4).
Newton's method consistently fails to find the correct optimum in the hierarchical
setting due to Hessian ill-conditioning, whereas BFGS --- which builds a positive-definite
inverse-Hessian approximation from gradient information --- succeeds at a small fraction
of the Gibbs cost.

\printbibliography

\end{document}
